\MathBookSourcePage{231}
\subsection{引理~\ref{lem:ch08-weighted-green-error} 的证明}
\label{sec:ch08-proof-weighted-green-error}
\setcounter{thmcount}{0}
\setcounter{equation}{0}

先证明有限元误差的一般加权范数估计. 它与第~\ref{chap:ch00-basic-concepts} 章中的估计相似, 但由于权函数 $\sigma_z$ 的特殊性质, 这里的结论稍强.

\MBSourceItem{1}
\begin{Proposition}
\label{prop:ch08-weighted-fem-error}
设 $d=2$ 或 $3$, 且式~\eqref{eq:ch08-primal-regularity} 对某个 $\mu>d$ 成立. 设有限元空间满足式~\eqref{eq:ch08-weighted-interpolation-estimates}、式~\eqref{eq:ch08-interpolation-orthogonality} 和式~\eqref{eq:ch08-weighted-inverse-estimate}. 令 $w$ 解式~\eqref{eq:ch08-primal-problem}, $w_h$ 解式~\eqref{eq:ch08-galerkin-problem}. 对任意 $\lambda>0$ 和 $\kappa>1$, 存在与 $z\in\overline\Omega$ 及 $h$ 无关的 $C<\infty$, 使得
\[
\MathBookFormulaID{formula-ch08-weighted-fem-error}
\begin{aligned}
\int_\Omega\sigma_z^{d+\lambda}|\nabla(w-w_h)|^2\,dx
&\leq C\int_\Omega\sigma_z^{d+\lambda-2}(w-w_h)^2\,dx\\
&\quad+C\int_\Omega\left(\sigma_z^{d+\lambda-2}(w-I^hw)^2
+\sigma_z^{d+\lambda}|\nabla(w-I^hw)|^2\right)\,dx.
\end{aligned}
\]
对所有 $z\in\Omega$ 和所有 $h$ 成立.
\end{Proposition}

\begin{Proof}
令 $e=w-w_h$, $\widetilde e=I^hw-w_h$, 并置 $\psi=\sigma_z^{d+\lambda}\widetilde e$. 注意对所有 $v\in V_h$ 有 $a(e,v)=0$. 由式~\eqref{eq:ch08-pointwise-ellipticity},
\MBSourceItem{2}
\begin{align}
\MathBookFormulaID{formula-ch08-weighted-error-coercivity-expansion}
C_a\int_\Omega\sigma_z^{d+\lambda}|\nabla e|^2\,dx
&\leq\int_\Omega\sigma_z^{d+\lambda}\sum_{ij}a_{ij}e_{,i}e_{,j}\,dx\notag\\
&=a(e,\sigma_z^{d+\lambda}e)-\int_\Omega\sum_{ij}a_{ij}e_{,i}e(\sigma_z^{d+\lambda})_{,j}\,dx\notag\\
&\quad-\int_\Omega\left(\sum_i b_ie_{,i}\sigma_z^{d+\lambda}e+b_0\sigma_z^{d+\lambda}e^2\right)\,dx\notag\\
&=a(e,\sigma_z^{d+\lambda}(w-I^hw)+\psi)-\int_\Omega\sum_{ij}a_{ij}e_{,i}e(\sigma_z^{d+\lambda})_{,j}\,dx\notag\\
&\quad-\int_\Omega\left(\sum_i b_ie_{,i}\sigma_z^{d+\lambda}e+b_0\sigma_z^{d+\lambda}e^2\right)\,dx\notag\\
&=a(e,\sigma_z^{d+\lambda}(w-I^hw))+a(e,\psi-I^h\psi)\notag\\
&\quad-\int_\Omega\left(\sum_{ij}a_{ij}e_{,i}e(\sigma_z^{d+\lambda})_{,j}
+\sum_i b_ie_{,i}\sigma_z^{d+\lambda}e+b_0\sigma_z^{d+\lambda}e^2\right)\,dx.
\label{eq:ch08-weighted-error-coercivity-expansion}
\end{align}
这里 $v_{,j}$ 表示关于第 $j$ 个坐标的偏导数. 式~\eqref{eq:ch08-weighted-error-coercivity-expansion} 中的最后两项可估计为
\begin{align*}
\MathBookFormulaID{formula-ch08-weighted-error-lower-order-estimate}
&\left|\int_\Omega\left(\sum_{ij}a_{ij}e_{,i}e(\sigma_z^{d+\lambda})_{,j}
+\sum_i b_ie_{,i}\sigma_z^{d+\lambda}e+b_0\sigma_z^{d+\lambda}e^2\right)\,dx\right|\\
&\quad\leq C\int_\Omega\bigl(|\nabla e||e|\sigma_z^{d+\lambda-1}
+\sigma_z^{d+\lambda}e^2\bigr)\,dx
\qquad\text{(由式~\eqref{eq:ch08-weight-infinity-bound} 和式~\eqref{eq:ch08-weight-derivative-bound})}\\
&\quad\leq C\left(\int_\Omega\sigma_z^{d+\lambda}|\nabla e|^2\,dx\right)^{1/2}
\left(\int_\Omega\sigma_z^{d+\lambda-2}e^2\,dx\right)^{1/2}
+C\int_\Omega\sigma_z^{d+\lambda}e^2\,dx\\
&\quad\leq\frac{C_a}{4}\int_\Omega\sigma_z^{d+\lambda}|\nabla e|^2\,dx
+C\int_\Omega\sigma_z^{d+\lambda-2}e^2\,dx.
\end{align*}
这里依次使用了 Schwarz 不等式、式~\eqref{eq:ch08-weight-infinity-bound} 和 Young 不等式~\eqref{eq:ch00-young-inequality-delta}. 同样, 第一项满足
\begin{align*}
\MathBookFormulaID{formula-ch08-weighted-error-approximation-estimate}
|a(e,\sigma_z^{d+\lambda}(w-I^hw))|
&\leq C\int_\Omega|\nabla e|
\bigl(\sigma_z^{d+\lambda}|\nabla(w-I^hw)|
+\sigma_z^{d+\lambda-1}|w-I^hw|\bigr)\,dx\\
&\quad+C\int_\Omega|e|\sigma_z^{d+\lambda}|w-I^hw|\,dx\\
&\leq C\left(\int_\Omega\sigma_z^{d+\lambda}(|\nabla e|^2+e^2)\,dx\right)^{1/2}\\
&\quad\times\left(\int_\Omega\sigma_z^{d+\lambda-2}(w-I^hw)^2
+\sigma_z^{d+\lambda}|\nabla(w-I^hw)|^2\,dx\right)^{1/2}\\
&\leq\frac{C_a}{4}\int_\Omega\sigma_z^{d+\lambda}(|\nabla e|^2+e^2)\,dx\\
&\quad+C\int_\Omega\sigma_z^{d+\lambda-2}(w-I^hw)^2
+\sigma_z^{d+\lambda}|\nabla(w-I^hw)|^2\,dx.
\end{align*}
第二项也可类似估计:
\begin{align*}
\MathBookFormulaID{formula-ch08-weighted-error-superapproximation-estimate}
|a(e,\psi-I^h\psi)|
&\leq C\int_\Omega\bigl(|\nabla e|(|\nabla(\psi-I^h\psi)|+|\psi-I^h\psi|)
+|e||\psi-I^h\psi|\bigr)\,dx\\
&\leq C\left(\int_\Omega\sigma_z^{d+\lambda}(|\nabla e|^2+e^2)\,dx\right)^{1/2}\\
&\quad\times\left(\int_\Omega\sigma_z^{-d-\lambda}
\bigl(|\nabla(\psi-I^h\psi)|^2+(\psi-I^h\psi)^2\bigr)\,dx\right)^{1/2}\\
&\leq\frac{C_a}{4}\int_\Omega\sigma_z^{d+\lambda}(|\nabla e|^2+e^2)\,dx\\
&\quad+C\int_\Omega\sigma_z^{-d-\lambda}
\bigl(|\nabla(\psi-I^h\psi)|^2+(\psi-I^h\psi)^2\bigr)\,dx.
\end{align*}
将这些估计代入式~\eqref{eq:ch08-weighted-error-coercivity-expansion}, 并使用式~\eqref{eq:ch08-weight-infinity-bound}, 得到
\MBSourceItem{3}
\begin{equation}
\MathBookFormulaID{formula-ch08-weighted-error-absorbed-bound}
\label{eq:ch08-weighted-error-absorbed-bound}
\begin{aligned}
\frac{C_a}{4}\int_\Omega\sigma_z^{d+\lambda}|\nabla e|^2\,dx
&\leq C\int_\Omega\sigma_z^{d+\lambda-2}e^2\,dx\\
&\quad+\int_\Omega\sigma_z^{-d-\lambda}
\bigl(|\nabla(\psi-I^h\psi)|^2+(\psi-I^h\psi)^2\bigr)\,dx
\\
&\quad+\int_\Omega\left(\sigma_z^{d+\lambda-2}(w-I^hw)^2
+\sigma_z^{d+\lambda}|\nabla(w-I^hw)|^2\right)\,dx.
\end{aligned}
\end{equation}
再由式~\eqref{eq:ch08-weighted-interpolation-estimates}、式~\eqref{eq:ch08-interpolation-orthogonality}、式~\eqref{eq:ch08-weight-derivative-bound} 和式~\eqref{eq:ch08-weighted-inverse-estimate},
\begin{align*}
\MathBookFormulaID{formula-ch08-superapproximation-full-chain}
&\int_\Omega\sigma_z^{-d-\lambda}
\bigl(|\nabla(\psi-I^h\psi)|^2+(\psi-I^h\psi)^2\bigr)\,dx\\
&\quad\leq\sum_{r=k}^{k_{\max}}h^{2r-2}
\int_\Omega\sigma_z^{-d-\lambda}|\nabla_{\mathcal E,r}\psi|^2\,dx\\
&\quad\leq C\sum_{r=k}^{k_{\max}}h^{2r-2}
\sum_{j=0}^{r-1}\int_\Omega\sigma_z^{d+\lambda-2(r-j)}|\nabla_j\widetilde e|^2\,dx\\
&\quad\leq C\sum_{r=k}^{k_{\max}}\sum_{j=0}^{r-1}h^{2r-2j-2}
\int_\Omega\sigma_z^{d+\lambda-2(r-j)}\widetilde e^2\,dx\\
&\quad\leq C\int_\Omega\sigma_z^{d+\lambda-2}\widetilde e^2\,dx.
\end{align*}
代回式~\eqref{eq:ch08-weighted-error-absorbed-bound} 即得命题结论. 还可将前一估计写成
\MBSourceItem{4}
\begin{equation}
\MathBookFormulaID{formula-ch08-superapproximation-estimate}
\label{eq:ch08-superapproximation-estimate}
\int_\Omega\sigma_z^{-d-\lambda}|\nabla(\psi-I^h\psi)|^2\,dx
\leq C\int_\Omega\sigma_z^{-d-\lambda-2}\psi^2\,dx.
\end{equation}
这种估计通常称为超逼近估计(Nitsche and Schatz 1974), 因为右端含较弱的范数, 且没有 $h$ 的负幂.
\end{Proof}

为估计命题~\ref{prop:ch08-weighted-fem-error} 右端含 $(w-w_h)^2$ 的项, 使用如下结论.

\MBSourceItem{5}
\begin{Proposition}
\label{prop:ch08-weighted-l2-error-absorption}
设 $d=2$ 或 $3$, 式~\eqref{eq:ch08-primal-regularity} 对某个 $\mu>d$ 成立, 且 $0<\lambda<2(1-d/\mu)$. 设有限元空间满足式~\eqref{eq:ch08-weighted-interpolation-estimates}、式~\eqref{eq:ch08-interpolation-orthogonality} 和式~\eqref{eq:ch08-weighted-inverse-estimate}. 令 $w$ 解式~\eqref{eq:ch08-primal-problem}, $w_h$ 解式~\eqref{eq:ch08-galerkin-problem}. 则存在 $\kappa_1<\infty$, 使得对所有 $\kappa\geq\kappa_1$ 和所有 $h$,
\[
\MathBookFormulaID{formula-ch08-weighted-l2-error-absorption}
\int_\Omega\sigma_z^{d+\lambda-2}(w-w_h)^2\,dx
\leq\varepsilon\int_\Omega\sigma_z^{d+\lambda}|\nabla(w-w_h)|^2\,dx.
\]
\end{Proposition}

\begin{Proof}
使用对偶论证. 令 $v\in V$ 解
\MBSourceItem{6}
\begin{equation}
\MathBookFormulaID{formula-ch08-dual-weighted-error-problem}
\label{eq:ch08-dual-weighted-error-problem}
a(\phi,v)=(\sigma_z^{d+\lambda-2}e,\phi)
\qquad\forall\phi\in V.
\end{equation}
于是, 利用 Schwarz 不等式、Young 不等式~\eqref{eq:ch00-young-inequality-delta} 和式~\eqref{eq:ch08-weighted-interpolation-estimates},
\begin{align*}
\MathBookFormulaID{formula-ch08-dual-weighted-error-chain}
\int_\Omega\sigma_z^{d+\lambda-2}e^2\,dx
&=a(e,v)=a(e,v-I^hv)\\
&\leq C\left(\int_\Omega\sigma_z^{d+\lambda}(|\nabla e|^2+e^2)\,dx\right)^{1/2}\\
&\quad\times\left(\int_\Omega\sigma_z^{-d-\lambda}
\bigl(|\nabla(v-I^hv)|^2+(v-I^hv)^2\bigr)\,dx\right)^{1/2}\\
&\leq\varepsilon\int_\Omega\sigma_z^{d+\lambda}(|\nabla e|^2+e^2)\,dx\\
&\quad+\frac{C}{\varepsilon}\int_\Omega\sigma_z^{-d-\lambda}
\bigl(|\nabla(v-I^hv)|^2+(v-I^hv)^2\bigr)\,dx\\
&\leq\varepsilon\int_\Omega\sigma_z^{d+\lambda}(|\nabla e|^2+e^2)\,dx
+\frac{Ch^2}{\varepsilon}\int_\Omega\sigma_z^{-d-\lambda}|\nabla_2v|^2\,dx.
\end{align*}
下一节将证明如下引理.

\MBSourceItem{7}
\begin{Lemma}
\label{lem:ch08-weighted-adjoint-regularity}
设 $d=2$ 或 $3$, 式~\eqref{eq:ch08-primal-regularity} 对某个 $\mu>d$ 成立, 且 $\lambda<2(1-d/\mu)$. 则方程
\[
\MathBookFormulaID{formula-ch08-adjoint-right-hand-side}
a(\phi,v)=(f,\phi)
\qquad\forall\phi\in V
\]
的解满足
\[
\MathBookFormulaID{formula-ch08-weighted-adjoint-regularity}
\int_\Omega\sigma^{-d-\lambda}|\nabla_2v|^2\,dx
\leq C\lambda^{-1}\zeta^{-2}
\int_\Omega\sigma^{4-d-\lambda}|\nabla f|^2\,dx,
\]
其中 $\sigma(x)=(|x-x^0|^2+\zeta^2)^{1/2}$, $x^0\in\overline\Omega$ 任意.
\end{Lemma}

将引理~\ref{lem:ch08-weighted-adjoint-regularity} 用于 $f=\sigma_z^{d+\lambda-2}e$, 得
\begin{align*}
\MathBookFormulaID{formula-ch08-dual-regularity-applied}
\int_\Omega\sigma_z^{-d-\lambda}|\nabla_2v|^2\,dx
&\leq C\lambda^{-1}(\kappa h)^{-2}
\int_\Omega\sigma_z^{4-d-\lambda}
\left|\nabla\bigl(\sigma_z^{d+\lambda-2}e\bigr)\right|^2\,dx\\
&\leq C\lambda^{-1}(\kappa h)^{-2}
\left(\int_\Omega\sigma_z^{d+\lambda}|\nabla e|^2\,dx
+\int_\Omega\sigma_z^{d+\lambda-2}e^2\,dx\right).
\end{align*}
因此
\[
\MathBookFormulaID{formula-ch08-weighted-l2-preabsorption}
\int_\Omega\sigma_z^{d+\lambda-2}e^2\,dx
\leq\varepsilon\int_\Omega\sigma_z^{d+\lambda}(|\nabla e|^2+e^2)\,dx
+\frac{C}{\lambda\kappa^2\varepsilon}
\left(\int_\Omega\sigma_z^{d+\lambda}|\nabla e|^2\,dx
+\int_\Omega\sigma_z^{d+\lambda-2}e^2\,dx\right).
\]
固定 $\varepsilon$ 和 $\lambda$, 取 $\kappa_1$ 足够大, 使得
\MBSourceItem{8}
\begin{equation}
\MathBookFormulaID{formula-ch08-weighted-l2-absorption-choice}
\label{eq:ch08-weighted-l2-absorption-choice}
\int_\Omega\sigma_z^{d+\lambda-2}e^2\,dx
\leq2\varepsilon\int_\Omega\sigma_z^{d+\lambda}|\nabla e|^2\,dx
\end{equation}
对所有 $\kappa\geq\kappa_1$ 成立. 重新命名 $\varepsilon$ 即得式中结论.
\end{Proof}

合并上述两个命题可见, 到有限元空间的 Galerkin 投影在某些加权空间中是拟最优的.

\MBSourceItem{9}
\begin{Theorem}
\label{thm:ch08-weighted-galerkin-quasioptimality}
设 $d=2$ 或 $3$, 式~\eqref{eq:ch08-primal-regularity} 对某个 $\mu>d$ 成立, 且 $0<\lambda<2(1-d/\mu)$. 设有限元空间满足式~\eqref{eq:ch08-weighted-interpolation-estimates}、式~\eqref{eq:ch08-interpolation-orthogonality} 和式~\eqref{eq:ch08-weighted-inverse-estimate}. 令 $w$ 解式~\eqref{eq:ch08-primal-problem}, $w_h$ 解式~\eqref{eq:ch08-galerkin-problem}. 则对所有 $z\in\overline\Omega$,
\[
\MathBookFormulaID{formula-ch08-weighted-galerkin-quasioptimality}
\int_\Omega\sigma_z^{d+\lambda}|\nabla(w-w_h)|^2+\sigma_z^{d+\lambda-2}(w-w_h)^2\,dx
\leq C\int_\Omega\sigma_z^{d+\lambda-2}(w-I^hw)^2
+\sigma_z^{d+\lambda}|\nabla(w-I^hw)|^2\,dx.
\]
\end{Theorem}

为完成引理~\ref{lem:ch08-weighted-green-error} 的证明, 将定理~\ref{thm:ch08-weighted-galerkin-quasioptimality} 应用于 $w=g^z$. 需证明
\MBSourceItem{10}
\begin{equation}
\MathBookFormulaID{formula-ch08-weighted-green-second-derivative}
\label{eq:ch08-weighted-green-second-derivative}
\int_\Omega\sigma_z^{d+\lambda}|\nabla_2g^z|^2\,dx\leq Ch^{\lambda-2}.
\end{equation}
这是下面先验估计的推论.

\MBSourceItem{11}
\begin{Lemma}
\label{lem:ch08-weighted-primal-regularity}
设 $d\leq3$, $0<\lambda<1$, 且当 $d=2$ 时 $a_{ij}\in W_p^1(\Omega)$ 对某个 $p>2$, 当 $d=3$ 时对某个 $p\geq12/5$; $i,j=1,\ldots,d$. 再假设式~\eqref{eq:ch08-primal-regularity} 对某个 $\mu>2d/(4-\lambda)$ 成立. 则
\[
\MathBookFormulaID{formula-ch08-weighted-primal-regularity}
a(\phi,v)=(\nu\cdot\nabla f,\phi)
\qquad\forall\phi\in V
\]
的解满足
\[
\MathBookFormulaID{formula-ch08-weighted-primal-regularity-bound}
\int_\Omega\sigma^{d+\lambda}|\nabla_2v|^2\,dx
\leq C\int_\Omega\sigma^{d+\lambda}|\nabla f|^2\,dx
+C\zeta^{-2}\int_\Omega\sigma^{d+\lambda}f^2\,dx,
\]
其中 $\sigma$ 为引理~\ref{lem:ch08-weighted-adjoint-regularity} 中定义的权函数.
\end{Lemma}
